% ── 问题一模板 ──
% 本文件负责 5.1 标题、问题一流程图和 5.1.x 分节装配。
% 流程图不设置 \subsubsection，直接位于“问题一的模型的建立和求解”标题之后。
% 问题二至题面实际第 n 问复制本文件后，统一修改标题、标签、文字和 \input 路径。

\subsection{问题一的模型的建立和求解}

\begin{figure}[H]
  \centering
  \begin{tikzpicture}[node distance=0.7cm and 0.55cm]
    \node[flowbox] (n11) {数据理解};
    \node[flowbox, right=of n11] (n12) {数据预处理};
    \node[flowbox, right=of n12] (n13) {变量关系};
    \node[flowbox, below=of n11] (n21) {模型假设};
    \node[flowbox, below=of n12] (n22) {模型建立};
    \node[flowbox, below=of n13] (n23) {参数求解};
    \node[flowbox, below=of n21] (n31) {模型检验};
    \node[flowbox, below=of n22] (n32) {结果分析};
    \node[flowbox, below=of n23] (n33) {结论输出};
    \draw[arrow] (n11) -- (n12);
    \draw[arrow] (n12) -- (n13);
    \draw[arrow] (n13) -- (n23);
    \draw[arrow] (n23) -- (n22);
    \draw[arrow] (n22) -- (n21);
    \draw[arrow] (n21) -- (n31);
    \draw[arrow] (n31) -- (n32);
    \draw[arrow] (n32) -- (n33);
  \end{tikzpicture}
  \caption{问题一建模思路图}
  \label{fig:problem1_flow}
\end{figure}

% ── 问题一：5.1.1 模型准备 ──
% 问题分析已迁移至第2章；问题流程图直接位于 5.1 标题下，不在本文件重复。
% 本文件只承载模型准备。问题二至题面实际第 n 问复制本文件并改编号、标签和问题文字。
% 数据预处理（仅问题一）五段：原因 → 数据理解与提取 → 处理步骤（含公式）→ 前后统计量对比 → 总结。
% 算法选择五段：问题本质定性 → 候选算法与评估维度 → 对比表（写真实算法名与等级）→ 分析对比结论 → 总结引出。

\subsubsection{模型准备}

在进行建模之前，需要……。本节将从数据预处理和算法选择两方面进行准备。

原始数据来源于……，包含……条记录和……个字段。经检查发现以下问题：……（缺失值、量纲不统一、异常值或编码问题）。

针对上述问题，按以下步骤进行预处理：第一步，……（缺失值处理及参数）；第二步，……（标准化或归一化，并写出公式）；第三步，……（特征编码或衍生）。

预处理后，数据规模为……行×……列，各指标统计量如下表所示。

\begin{table}[H]
  \centering
  \caption{数据预处理前后统计量对比}
  \begin{tabularx}{\textwidth}{CCCCC}
    \toprule
    指标 & 预处理前均值 & 预处理前标准差 & 预处理后均值 & 预处理后标准差 \\
    \midrule
    …… & …… & …… & …… & …… \\
    \bottomrule
  \end{tabularx}
  \label{tab:preprocess_compare1}
\end{table}

综上所述，数据预处理完成，为后续建模提供了清洁的输入数据。

该问题本质上是一个……问题，需要在……条件下……。常用求解算法包括 A 算法、B 算法和 C 算法等。本节从……、……和……三方面进行评估。

\begin{table}[H]
  \centering
  \caption{算法能力对比}
  \begin{tabularx}{\textwidth}{>{\hsize=2.5\hsize\centering\arraybackslash}X>{\hsize=0.6\hsize\centering\arraybackslash}X>{\hsize=0.6\hsize\centering\arraybackslash}X>{\hsize=0.6\hsize\centering\arraybackslash}X>{\hsize=0.7\hsize\centering\arraybackslash}X}
    \toprule
    算法名称 & 维度1 & 维度2 & 维度3 & 稳定性 \\
    \midrule
    算法 A（缩写） & 优 & 良 & 中 & 良 \\
    算法 B（缩写） & 良 & 优 & 良 & 优 \\
    算法 C（缩写） & 中 & 中 & 优 & 中 \\
    \bottomrule
  \end{tabularx}
  \label{tab:algorithm_compare1}
\end{table}

通过对比可以看出，算法 A 在……方面表现最优，……。因此本文选择算法 A 作为核心求解方法。

综上所述，本节完成了数据预处理和算法选型，接下来将建立数学模型并进行求解。


\input{5.1.2.建模与求解.tex}
